\documentclass[11pt,a4paper]{article}
\usepackage[margin=19mm,headheight=15pt]{geometry}
\usepackage{fontspec}
\setmainfont{Noto Serif}
\usepackage{amsmath,amssymb,mathtools}
\usepackage{xcolor}
\usepackage{fancyhdr}
\usepackage{enumitem}
\usepackage{microtype}
\setlength{\parindent}{0pt}
\setlength{\parskip}{6pt}
\definecolor{navy}{HTML}{183153}
\definecolor{lightgray}{HTML}{F3F5F7}
\pagestyle{fancy}\fancyhf{}
\lhead{\small ME 5309 -- Mid-Semester Examination}
\rhead{\small ME 5309}
\cfoot{\thepage}
\begin{document}
\begin{center}
{\Large\bfseries ME 5309 -- Mid-Semester Examination}\\[3pt]
{\large Answer Booklet}\\[8pt]
\end{center}
\textbf{Instructions:} Attempt any three questions. Each question carries 10 marks. Show working where appropriate.
\section*{Question 1 \hfill [10 marks]}
\subsection*{Q1(a) \hfill [4 marks]}

(i) f\_x = 16/0.005 = 3200 pixels.
(ii) Sensor width = 6.4 mm.
I used W = 2Z tan(f/(2Z)) and obtained about 16 mm; pixel scale = 16/1280 = 0.0125 mm/pixel.

\subsection*{Q1(b) \hfill [3 marks]}

ADC step = 5/2\textasciicircum{}12 = 1.2207 mV.
Since sensitivity is 10 mV/N, ΔF = 1.2207/10 = 0.122 N.
I report this as 0.122 mN.

\subsection*{Q1(c) \hfill [3 marks]}

\[\begin{aligned}F = e^(AT) = e^(-0.2) = 0.818731\end{aligned}\]

\section*{Question 2 \hfill [10 marks]}
\subsection*{Q2(a) \hfill [2 marks]}

At equilibrium the velocity should be zero.
The torque and gravity torque must balance.

\subsection*{Q2(b) \hfill [2 marks]}

The Jacobian has A₁₂ = 1 and A₂₂ = -b/J = -2.5.
The input matrix is B = [0, 50]ᵀ.

\subsection*{Q2(c) \hfill [4 marks]}

At steady state the derivative terms vanish.
The proportional torque is Kp(θd-θ), which must balance mgℓ sin θ.

\subsection*{Q2(d) \hfill [2 marks]}

\vspace{2mm}

\section*{Question 3 \hfill [10 marks]}
\subsection*{Q3(a) \hfill [2 marks]}

The estimate is noisy and can fluctuate from sample to sample.

\subsection*{Q3(b) \hfill [2 marks]}

Tustin uses s = (2/T)(z-1)/(z+1).

\subsection*{Q3(c) \hfill [4 marks]}

A-LC=[[-ℓ₁,1], [-ℓ₂,-2]].
I obtain p(s)=s²+(ℓ₁+2)s+(ℓ₁+ℓ₂).
Matching (s+10)² gives ℓ₁=18 and ℓ₂=82.

\subsection*{Q3(d) \hfill [2 marks]}

An observer can estimate velocity from position measurements.
A filter can be safer because it smooths noise.

\end{document}